% Section 2: Background and Preliminaries
\section{Background and Preliminaries}
\label{sec:bg}

\subsection{Reinforcement Learning Formulation for LLMs}

In the context of LLM alignment, the training process is formulated as a contextual bandit problem rather than a full Markov Decision Process (MDP). This formulation is appropriate because the model generates a complete response autoregressively before receiving any reward signal. Formally, we define:

\begin{itemize}[leftmargin=*, itemsep=1pt]
\item \textbf{State/Context} $x \sim \mathcal{D}$: The input prompt sampled from the training distribution $\mathcal{D}$.
\item \textbf{Action/Response} $y = (y_1, y_2, \ldots, y_T)$: The token sequence generated by the model, where $T$ is the response length.
\item \textbf{Policy} $\pi_\theta(y|x) = \prod_{t=1}^{T} \pi_\theta(y_t | x, y_{<t})$: The autoregressive language model parameterized by $\theta$.
\item \textbf{Reward} $r(x, y) \in \mathbb{R}$: A scalar signal evaluating the quality of response $y$ given prompt $x$. This may come from a trained reward model, a rule-based verifier, or human annotation.
\end{itemize}

The optimization objective is to find policy parameters $\theta^*$ that maximize the expected reward:
\begin{equation}
\theta^* = \arg\max_\theta J(\theta) = \arg\max_\theta \E_{x \sim \mathcal{D},\ y \sim \pi_\theta(\cdot|x)}\left[r(x, y)\right]
\label{eq:obj}
\end{equation}

A critical challenge is preventing the policy from deviating too far from the reference model $\pi_{\text{ref}}$ (typically the SFT checkpoint), as unconstrained optimization can lead to reward hacking and mode collapse. Different algorithms address this constraint through different mechanisms: explicit KL penalties (GRPO), implicit regularization (DPO), clipped surrogate objectives (PPO, DAPO), or combinations thereof.

\subsection{The RLHF Pipeline in Detail}

The standard RLHF pipeline consists of three stages:

\textbf{Stage 1 --- Supervised Fine-Tuning (SFT):} A pre-trained language model is fine-tuned on high-quality demonstration data $\{(x_i, y_i^*)\}_{i=1}^{N}$ using maximum likelihood estimation:
\begin{equation}
\mathcal{L}_{\text{SFT}}(\theta) = -\E_{(x,y^*) \sim \mathcal{D}_{\text{demo}}}\left[\log \pi_\theta(y^*|x)\right]
\end{equation}
The resulting model $\pi_{\text{SFT}}$ serves as the initialization for subsequent RL training and as the reference model $\pi_{\text{ref}}$.

\textbf{Stage 2 --- Reward Model Training:} A reward model $r_\phi(x, y)$ is trained on human preference pairs $(x, y_w, y_l)$ where $y_w \succ y_l$ denotes that $y_w$ is preferred over $y_l$. The Bradley-Terry model provides the training objective:
\begin{equation}
\mathcal{L}_{\text{RM}}(\phi) = -\E\left[\log \sigma\left(r_\phi(x, y_w) - r_\phi(x, y_l)\right)\right]
\end{equation}
where $\sigma(\cdot)$ is the logistic sigmoid function.

\textbf{Stage 3 --- RL Policy Optimization:} The language model policy $\pi_\theta$ is optimized against the reward model using an RL algorithm. This is the stage where the choice of algorithm (PPO, GRPO, DAPO) matters most. DPO bypasses stages 2 and 3 entirely by directly optimizing on preference pairs.

\subsection{Unified Notation}

Table~\ref{tab:notation} establishes the unified notation used throughout this survey to enable direct cross-algorithm comparison.

\begin{table}[ht]
\centering
\caption{Unified Mathematical Notation Used in This Survey}
\label{tab:notation}
\begin{tabular}{cl}
\toprule
\textbf{Symbol} & \textbf{Definition} \\
\midrule
$\pi_\theta$ & Current policy network (LLM being trained) \\
$\pi_{\text{ref}}$ & Reference (frozen) policy, typically SFT model \\
$\pi_{\theta_{\text{old}}}$ & Policy parameters from previous update step \\
$r(x,y)$ & Reward function (learned or rule-based) \\
$\hat{A}$ & Estimated advantage function \\
$r_t(\theta)$ & Importance sampling ratio $\pi_\theta/\pi_{\theta_{\text{old}}}$ \\
$\varepsilon$ & PPO/GRPO symmetric clipping parameter \\
$\varepsilon_{\text{low}}, \varepsilon_{\text{high}}$ & DAPO asymmetric clipping bounds \\
$\beta$ & KL penalty coefficient or DPO temperature \\
$G$ & Group size for GRPO/DAPO sampling \\
$\gamma, \lambda$ & Discount factor and GAE decay parameter \\
$V_\phi$ & Value network (Critic) for PPO \\
$y_w, y_l$ & Preferred and rejected responses (DPO) \\
$|o_i|$ & Token length of the $i$-th response \\
\bottomrule
\end{tabular}
\end{table}
